\documentclass[10pt,a4paper]{amsart}
\usepackage[margin=19mm]{geometry}
\usepackage{amsmath,amssymb,mathtools}
\usepackage[scr=rsfs,cal=euler]{mathalfa}
\usepackage{xfrac}
\usepackage[unicode,hidelinks,bookmarks=false]{hyperref}
\newcommand{\ie}{i.e.\ }
\newcommand{\cf}{cf.\ }
\newcommand{\resp}{respectively\ }
\newcommand{\Subsection}{\subsection}
\providecommand{\nobreakdash}{\nobreak\hbox{-}}
\setlength{\emergencystretch}{2em}
\multlinegap0pt
\usepackage{mathtools}
\usepackage{amssymb}
\usepackage{enumerate}
\usepackage{xcolor}
\newtheorem{theorem}{Theorem}
\newcommand{\Ann}{\operatorname{Ann}\nolimits}
\newcommand{\Ext}{\operatorname{Ext}\nolimits}
\newcommand{\Hom}{\operatorname{Hom}\nolimits}
\newcommand{\add}{\mathsf{add}\hspace{.01in}}
\newcommand{\perpall}[1]{{}^{\perp_{\geq 0}}#1}
\title{Self-Orthogonal Tau-Tilting Modules and Tilting Modules II: Annihilator Separation}
\author{Xiaojin Zhang}
\date{Source: arXiv:2608.27937v1 (CC BY 4.0)}
\begin{document}
\maketitle

\noindent\emph{Source and licence.} Self-Orthogonal Tau-Tilting Modules and Tilting Modules II: Annihilator Separation. Version arXiv:2608.27937v1 (CC BY 4.0), distributed by arXiv under the Creative Commons Attribution 4.0 International licence, http://creativecommons.org/licenses/by/4.0/, as recorded in the arXiv licence metadata for this exact version and shown on the abstract page https://arxiv.org/abs/2608.27937v1. Full licence notice: \texttt{licenses/arxiv-2608.27937-CC-BY-4.0.txt}.\par\smallskip

\noindent\textit{Editorial scope.} Counted unit: the article's theorem thm:orthogonality (source0.tex lines 306--316). The article's complete original proof of it is delivered with it (source0.tex lines 318--345, one \texttt{proof} environment of 28 lines). No mathematics is written, repaired or re-derived: the statement and the proof are the article's own lines, in the article's own notation, and no retained line is substituted or re-flowed.  Retained uncounted as the source-local context the proof needs: lines 276--279.  Nothing was omitted from any retained span, so the package carries no omission record.  No retained line contains a citation, so the wrapper carries no bibliography and no bibliography entry.  No retained line contains a figure environment, an image-inclusion command or drawing code, so no image is delivered and the package declares no asset dependency.  Editorial formatting only, in addition: an AMS \texttt{amsart} preamble that restates the article's own declarations and macros used by the retained lines, namely the environments \texttt{theorem} on separate counters, together with the standard packages those lines need.  The retained environments print in this document as Theorem 1 in order of appearance, whereas the article prints the same items with the numbering of its own full document.  The complete native source of the article as distributed in the arXiv e-print for this exact version is bundled unchanged as \texttt{sources/208749/source0.tex}.
\begin{equation}\label{eq:canonical-approximation}
 A\xrightarrow{f}T_0\longrightarrow T_1\longrightarrow0,
 \qquad T_0,T_1\in\add T,
\end{equation}

\begin{theorem}[Complete annihilator orthogonality]
\label{thm:orthogonality}
Let $T$ be a self-orthogonal $\tau$-tilting $A$-module and set
$I=\Ann_A(T)$.  Then
\[
              \Hom_A(I,T)=0
              \quad\text{and}\quad
              \Ext_A^n(I,T)=0\quad(n\geq1).
\]
Equivalently, $I\in\perpall{T}$.
\end{theorem}

\begin{proof}
Choose \eqref{eq:canonical-approximation} and put
$C=\operatorname{Im}f$.  It gives short exact sequences
\[
 0\longrightarrow C\longrightarrow T_0\longrightarrow T_1
 \longrightarrow0
\]
and
\[
 0\longrightarrow I\longrightarrow A\longrightarrow C
 \longrightarrow0.
\]
Because $T_0,T_1\in\add T$ and $T$ is self-orthogonal, the first
sequence yields
\[
                    \Ext_A^n(C,T)=0\qquad(n\geq1).
\]
The second sequence then gives $\Ext_A^n(I,T)=0$ for every $n\geq1$.

Applying $\Hom_A(-,T)$ to the second sequence also shows that the
restriction map
\[
                 \Hom_A(A,T)\longrightarrow\Hom_A(I,T)
\]
is surjective. Every map $A\to T$ factors through the left
$\add T$-approximation $f$ and therefore vanishes on $I=\ker f$.
The restriction map is thus zero, and $\Hom_A(I,T)=0$.
\end{proof}

\end{document}
